% !TeX root = ../本科微分几何笔记.tex
% 由主文档通过 \input 载入；请编译主文档。
\setcounter{chapter}{4} % 保留原书章号。
\chapter{第一基本形式与内蕴度量}
\label{chap:dg-metric}

\section{由切向量内积到第一基本形式}
\begin{definition}[第一基本形式]
曲面在 $p$ 处的\emph{第一基本形式（first fundamental form）}是欧氏内积在 $T_p\Sigma$ 上的限制，记为 $\mathrm I_p(X,Y)=\ip XY$。在坐标 $(u^1,u^2)=(u,v)$ 下，记 $\mathbf x_i=\partial_i\mathbf x$，并令
\[
 (g_{ij})=\begin{pmatrix}E&F\\F&G\end{pmatrix},\qquad
 E=\ip{\mathbf x_u}{\mathbf x_u},\quad
 F=\ip{\mathbf x_u}{\mathbf x_v},\quad
 G=\ip{\mathbf x_v}{\mathbf x_v}.
\]
常将其写成 $\mathrm I=E\dd u^2+2F\dd u\dd v+G\dd v^2$，此处是二次型的记号，不是对坐标微分进行普通数值相乘。
\end{definition}
正则性保证 $E>0$、$G>0$ 及 $EG-F^2>0$。对于坐标列 $a,b$，切向量的长度与夹角由
\[
 \norm X^2=a^{\mathsf T}\mathbf g a,\qquad
 \cos\angle(X,Y)=\frac{a^{\mathsf T}\mathbf g b}
 {\sqrt{a^{\mathsf T}\mathbf g a}\sqrt{b^{\mathsf T}\mathbf g b}}
\]
决定，其中角公式要求 $X,Y\ne0$。

\begin{proposition}[度量变换律]
若旧坐标 $u$ 与新坐标 $\widetilde u$ 满足 $u=\phi(\widetilde u)$，Jacobi 矩阵为 $J=D\phi$，则
\begin{equation}
 \widetilde{\mathbf g}=J^{\mathsf T}\mathbf gJ.
 \label{eq:dg-metric-transform}
\end{equation}
\end{proposition}
\begin{proof}
链式法则给出 $\widetilde{\mathbf x}_a=\sum_i\mathbf x_i\partial u^i/\partial\widetilde u^a$，两边作内积即得。向量坐标同时作逆向变换，因而计算出的内积保持不变。
\end{proof}

\section{曲线长度、面积与曲面距离}
若曲面曲线为 $\gamma(t)=\mathbf x(u(t),v(t))$，则
\begin{equation}
 L(\gamma)=\int_a^b\sqrt{E\dot u^2+2F\dot u\dot v+G\dot v^2}\dd t.
 \label{eq:dg-surface-length}
\end{equation}
局部小矩形经参数化后，一阶近似为由 $\mathbf x_u\dd u,\mathbf x_v\dd v$ 张成的平行四边形，故面积元为
\begin{equation}
 \dd A=\norm{\mathbf x_u\times\mathbf x_v}\dd u\dd v
       =\sqrt{EG-F^2}\dd u\dd v.
 \label{eq:dg-area}
\end{equation}
式~\eqref{eq:dg-metric-transform} 给出 $\det\widetilde{\mathbf g}=(\det J)^2\det\mathbf g$，配合二重积分的换元公式说明面积不依赖参数。面积元取正值，所以不需要整体定向；定向面积形式则还记录方向，二者应区分。

\begin{definition}[内蕴距离]
在连通曲面上，定义
\[
 d_\Sigma(p,q)=\inf\{L(\gamma):\gamma\text{ 为连接 }p,q\text{ 的分段光滑曲线}\}.
\]
连通曲面局部道路连通，从而道路连通，并可用坐标内的分段光滑路径连接两点，因此集合非空。
\end{definition}
长度不小于端点之间的欧氏距离，所以 $d_\Sigma(p,q)\ge\norm{p-q}$，这保证距离的正定性；对称性来自反向曲线，三角不等式来自拼接近乎达到下确界的曲线。下确界未必能由某条曲线达到，例如去掉原点的平面上连接 $(-1,0)$ 与 $(1,0)$ 的最短长度下确界为 $2$，却没有长度恰好为 $2$ 的允许路径。

\section{局部等距与保角映射}
\begin{definition}[等距与保角]
光滑映射 $F:\Sigma\to\widetilde\Sigma$ 若在每点满足
\[
 \widetilde{\mathrm I}_{F(p)}(\dd F_pX,\dd F_pY)=\mathrm I_p(X,Y),
\]
称为\emph{局部等距（local isometry）}。正定性迫使 $\dd F_p$ 可逆，故它也是局部微分同胚。若 $F$ 还是整体微分同胚，则称为等距。若上式右边改为 $\lambda(p)^2\mathrm I_p(X,Y)$，其中 $\lambda>0$，称为\emph{保角映射（conformal map）}。
\end{definition}
局部等距保持每条曲线的长度，但若不是整体双射，未必保持任意两点的整体距离。保角映射保持非零切向量之间的角度，长度则乘以局部因子 $\lambda$，面积乘以 $\lambda^2$。

\begin{example}[圆柱展开]\label{exam:dg-cylinder-isometry}
半径为 $R$ 的圆柱取参数 $\mathbf x(u,v)=(R\cos(u/R),R\sin(u/R),v)$。由于 $E=G=1$、$F=0$，映射从平面到圆柱处处局部等距。它在整个平面上不单射，因为 $u$ 与 $u+2\pi R$ 映到同一点；限制在宽度小于 $2\pi R$ 的开条带上才得到到其像的等距。因此“圆柱可展开”首先是局部长度保持的陈述。
\end{example}

\begin{example}[球面的立体投影坐标]\label{exam:dg-stereographic}
单位球面去掉北极的一个参数化为
\[
 \mathbf x(u,v)=\frac1{1+u^2+v^2}(2u,2v,u^2+v^2-1).
\]
直接求导作内积得
\begin{equation}
 \mathrm I=\frac4{(1+u^2+v^2)^2}(\dd u^2+\dd v^2).
 \label{eq:dg-stereo-metric}
\end{equation}
这说明立体投影坐标保角，却不是等距；离原点越远，参数平面上的同样长度在球面上对应越小的长度。不能把保角误解成保面积。
\end{example}

\section{习题}
\begin{exercise}\label{ex:dg-metric-graph}
对图像 $z=f(u,v)$ 求第一基本形式与面积元，并用于 $f(u,v)=uv$。
\end{exercise}
\begin{exercise}\label{ex:dg-metric-sphere}
求半径 $R$ 球面的经纬度量、纬线长度及整个球面的面积。
\end{exercise}
\begin{exercise}\label{ex:dg-metric-cylinder}
圆柱上两点的角坐标差为 $\Delta\theta$，高度差为 $\Delta z$。证明内蕴距离为 $\min_{m\in\mathbb Z}\sqrt{R^2(\Delta\theta+2\pi m)^2+\Delta z^2}$。
\end{exercise}
\begin{exercise}\label{ex:dg-metric-conformal}
设两个二维内积空间间的可逆线性映射保持所有非零向量之间的角度，证明它把内积乘以同一个正常数。
\end{exercise}

